\section{Robustness and Generalisation}\label{sec:robustness}
The deployed models are frozen and nothing here is tuned. These experiments re-use the test split \emph{after}
the one-time headline evaluation of Section~\ref{sec:test}, for evaluation only, with nothing fed back; each use
is recorded in the test-run log (two early runs were disclosed in a later note; Appendix~\ref{app:repro}).

\subsection{Robustness to Post-Processing}\label{sec:rob-degrade}
Every test image (1,892) was degraded \emph{before} the protocol was applied: JPEG re-compression at quality
95--50; downscaling by $\times$0.9, $\times$0.75 and $\times$0.5, with the ground-truth mask resized alongside;
Gaussian blur, $\sigma=0.5$--1.5; Gaussian noise, $\sigma=2$--10 grey levels; and a simulated social-media upload
($\times$0.8 downscale, then JPEG quality 70). Table~\ref{tab:rob-degrade} reports
the image-level AUC of the deployed forest and the pixel F1 of the deployed localiser (all conditions are plotted
in Fig.~\ref{fig:robustness}, Appendix~\ref{app:robustness}).

\begin{table}[t]
\centering
\caption{Image AUC (pixel F1) of the frozen detector on the degraded test set.}\label{tab:rob-degrade}
\footnotesize
\begin{tabular}{lcc}
\toprule
Condition & Protocol R & Protocol B \\
\midrule
Clean               & 0.805 (0.197) & 0.920 (0.308) \\
JPEG Q85            & 0.793 (0.193) & 0.903 (0.306) \\
JPEG Q75            & 0.799 (0.193) & \textbf{0.739} (0.223) \\
JPEG Q50            & 0.781 (0.191) & 0.719 (0.214) \\
Resize $\times$0.75 & 0.768 (0.174) & 0.747 (0.170) \\
Resize $\times$0.5  & \textbf{0.693} (0.126) & 0.698 (0.121) \\
Blur $\sigma=1.5$   & 0.764 (0.201) & 0.754 (0.203) \\
Noise $\sigma=10$   & 0.794 (0.197) & 0.803 (0.181) \\
Social upload       & \textbf{0.762} (0.176) & \textbf{0.721} (0.185) \\
\bottomrule
\end{tabular}
\end{table}


Protocol R is robust to re-compression and noise: its AUC stays within 0.03 of the clean value down to JPEG
quality 50 and noise $\sigma=10$, and its localisation F1 hardly changes. R already discards the grid-bound
compression traces, so its evidence (copy-move keypoints, texture-normalised ELA, edges) does not depend on them.
Strong downscaling ($\times$0.5) is its main weakness (AUC 0.69, F1 0.13): smaller images leave fewer keypoints
and blocks.

Protocol B's advantage is fragile. B survives a re-save at quality 95 or 85 (0.900, 0.903; 85 is B's own
re-encode quality), but at a lower quality its clean AUC (0.920) falls below protocol R's: 0.739 at Q75, 0.772 at
Q65 (the decline is not monotonic) and 0.719 at Q50. Under blur, noise or downscaling, B's advantage over R largely
disappears; B stays at most about 0.02 ahead (e.g.\ 0.820 vs 0.802 under noise $\sigma=2$). B's extra signal is
the within-image compression history (Section~\ref{sec:params}), and any further lossy processing overwrites it.
A simulated social-media upload costs about 0.04 AUC under R (0.805 $\to$ 0.762) and 0.20 under B
(0.920 $\to$ 0.721). For images that have been shared online the strict protocol is the realistic setting, which
supports choosing R as the main protocol, a choice Section~\ref{sec:classification} made on other grounds (its
pre-registered rule).

\subsection{Synthetic Forgeries with Known Masks}\label{sec:rob-synthetic}
From the test set's authentic images we generated 480 forgeries with exact masks (tampered regions cover a median
6\,\% of the image, range 0.2--22\,\%), 120 of each kind: hard-edged splices (a random region cut from another
authentic image), seamless splices (the same, blended in with Poisson cloning~\cite{perez2003poisson}), plain
copy-moves (a translated copy of a region) and rotated and scaled copy-moves ($\pm20^\circ$, $\times$0.85--1.15).
A further 240 untouched authentic images, never used as hosts or donors, serve as negatives. All images were
saved losslessly and evaluated as made and after a simulated social-media upload. Mask definitions are given in
Appendix~\ref{app:robustness}.

\begin{table*}[t]
\centering
\caption{Synthetic forgeries: AUC against the 240 untouched authentic images [95\,\% CI], share judged
``tampered'', and pixel F1 (for copy-moves also F1 against source $\cup$ copy, in parentheses). False alarms on the
untouched images: 1.2\,\% (R), 8.8\,\% (B).}\label{tab:rob-synthetic}
\footnotesize
\begin{tabular}{lccccc}
\toprule
Kind & R: AUC & R: detected & R: F1 (source $\cup$ copy) & B: AUC & B: F1 (source $\cup$ copy) \\
\midrule
Splice, hard edge          & 0.569\ci{0.508}{0.629} & 1\,\%  & 0.108 & 0.570 & 0.158 \\
Splice, seamless           & 0.552\ci{0.487}{0.609} & 2\,\%  & 0.079 & 0.551 & 0.225 \\
Copy-move, shift           & \textbf{0.885}\ci{0.830}{0.926} & 73\,\% & 0.376 (\textbf{0.597}) & 0.937 & 0.497 (0.703) \\
Copy-move, rotate + scale  & \textbf{0.896}\ci{0.857}{0.934} & 63\,\% & 0.360 (\textbf{0.571}) & 0.908 & 0.484 (0.653) \\
After a social upload (all kinds) & 0.693 & 27\,\% & 0.182 & 0.687 & 0.222 \\
\bottomrule
\end{tabular}
\end{table*}

Copy-moves are found, including rotated and scaled ones (Table~\ref{tab:rob-synthetic}: AUC 0.89--0.94, F1
0.36--0.50): mirror-aware, scale- and rotation-invariant keypoint matching does what it was designed for.
Fig.~\ref{fig:synthetic} shows that the localiser marks \emph{both} the source and the pasted copy. The ground
truth marks only the copy, so pixel F1 understates copy-move localisation; scored against source $\cup$ copy, F1
rises to 0.60/0.57 under R and 0.70/0.65 under B (shift / rotate + scale).

Splices between two CASIA authentic photos are essentially undetectable (AUC 0.55--0.57; the CI for seamless
splices includes 0.5). This is consistent with the test-set error analysis, where missed forgeries were mostly
splices (Section~\ref{sec:test}), while CASIA's own splices remain more detectable (AUC 0.68 under R). Here the
host and donor share the same kind of camera JPEG history, noise level and processing, so there is no compression,
noise or sharpness inconsistency for the modules to detect; CASIA's own splices are detectable to the extent that
their pasted content differs in such traces, and more so under B. Seamless blending removes the hard edge, yet even
the hard-edged splices are rarely flagged: edge sharpness alone is not enough evidence. Social-media processing
weakens detection (overall AUC 0.725 $\to$ 0.693 under R), with losses spread over the kinds: seamless splices
$-$0.051 (to chance, 0.501), rotated/scaled copy-moves $-$0.048, shifted copy-moves $-$0.024 and hard-edged
splices $-$0.005.

\begin{figure}[t]
\centering
\includegraphics[width=\columnwidth]{fig_synthetic_examples.png}
\caption{Synthetic forgeries from test-set authentic images, median-F1 example of each kind under protocol R
(rows: hard-edged splice, seamless splice, shifted copy-move, rotated and scaled copy-move). Columns: forged image,
true region, fused probability map, predicted mask. For copy-moves the localiser marks both the source and the
copy, whereas the ground truth marks only the copy; the splices are missed.}
\label{fig:synthetic}
\end{figure}

\subsection{Cross-Dataset Test: MICC-F220}\label{sec:rob-micc}
MICC-F220~\cite{amerini2011sift} (used for non-commercial research) contains 110 copy-move forgeries and 110
originals, with larger images than most of CASIA, and is labelled at image level only. The forgeries come from only
11 scenes, each with 10 attack variants (translation, rotation, scaling and combinations), so the AUC rests on 11
forged scenes; the other 99 originals have no forged version. In the bootstrap, each scene's original and its
forgeries form one group.

\begin{table}[t]
\centering
\caption{Cross-dataset evaluation on MICC-F220 without retraining (CM: copy-move keypoint matches).}\label{tab:rob-micc}
\footnotesize
\setlength{\tabcolsep}{4pt}
\begin{tabular}{lcc}
\toprule
 & Protocol R & Protocol B \\
\midrule
AUC [95\,\% CI]                 & \textbf{0.972\ci{0.933}{0.998}} & 0.902\ci{0.767}{0.984} \\
Share of images judged          & 48\,\% & 89\,\% \\
Accuracy on judged          & 0.943 & 0.790 \\
Forgeries judged tampered   & 86\,\% & 88\,\% \\
Originals judged tampered   & 5.5\,\% & 28\,\% \\
Forgeries with CM matches   & 97\,\% & 91\,\% \\
Originals with CM matches   & 23\,\% & 32\,\% \\
\bottomrule
\end{tabular}
\end{table}

The CASIA-trained detector transfers to this unseen copy-move dataset without retraining, with AUC 0.97 under R
(Table~\ref{tab:rob-micc}), driven by the keypoint copy-move module, which finds matches in 97\,\% of the
forgeries. The false-alarm mechanism of Section~\ref{sec:test} appears here too: 23\,\% of the originals contain
copy-move keypoint matches, presumably from repeated structure (not inspected). The calibrated model still keeps
false ``tampered'' verdicts at 5.5\,\% under R, mostly by abstaining: of the 110 originals, 100 fall in the
uncertain band, 4 are judged authentic and 6 tampered, so the judged accuracy of 0.943 comes almost entirely from
the forgeries. B generalises worse: its AUC is 0.90, and the paired group-bootstrap difference is
R $-$ B $=$ +0.069\ci{0.010}{0.172}. Its false-alarm share is 28\,\% against 5.5\,\% for R, although this
comparison uses each protocol's own uncertain band ($\pm0.04$ vs $\pm0.24$). This is consistent with B's reliance
on CASIA-specific compression history.

Other external sets, including COVERAGE~\cite{wen2016coverage}, could not be obtained automatically
(Appendix~\ref{app:robustness}), so a cross-dataset \emph{splicing} test is missing; given
Section~\ref{sec:rob-synthetic}, we expect splicing to be the weaker case.

In summary, under the strict protocol the detector keeps most of its accuracy under re-compression and noise,
generalises to an unseen copy-move dataset and is reliable for copy-move forgeries; its clear limit is splicing
between photographs with the same processing history. Protocol B looks stronger on clean CASIA images, but that
advantage vanishes after a re-save at JPEG quality 75 or below, and on an external dataset.
